\scanpage{61}
quitte à passer à une extension finie de $k_0$, $\sigma_k\rightsquigarrow({}^\sigma g)^{-1}g$ \inserted{au lieu de $g^{-1}g^\sigma$} où $g\in G(\unsure{k|K})$ est gén. $(k_0)$.
\margin{(\unsure{$g$ n'est|$g$ est} déterminé mod mult par $g_0\in G(k_0)$)}
\figure{\begin{tikzpicture}[scale=.5]
\draw (0,0) .. controls (1,.35) and (2.1,.2) .. (2.4,.3) .. controls (4,.15) and (6,.3) .. (7,.3);
\draw[->] (2.4,-.7)--(2.4,.3);
\draw (7.1,.6) .. controls (7.1,1.15) and (8.5,1) .. (9.1,1.25);
\end{tikzpicture}}
On trouve alors que $g^\sigma$ \uline{\unsure{$\forall$|tt}} pt gén. de $G(k)$, et par suite $g^\sigma g^{-1}$ prend \deleted{\unsure{tt}} \inserted{\unsure{exactt}} les valeurs $a$ telles que $\deleted{g}ag$ soit gén. $/k_0$, par ex. \unsure{tts|ttes} valeurs telles que \deleted{$(\unsure{ag})$} \deleted{$ag$} $a$ soit gén. sur $k_0(g)$.

\deleted{En d'autres termes} Considérons enfin l'hom total $\sigma\longrightarrow\sigma_*$
\diagram{\begin{tikzcd}[ampersand replacement=\&,column sep=large]
\operatorname{Gal}(k/k_0)\arrow[r,"{\deleted{\unsure{\varphi_\infty}}\inserted{\varphi_v}}"]\&G(\mathbb A_{\unsure{C|c}})
\end{tikzcd}}
\deleted{\unsure{de dire} On a envie de dire que l'image}

Considérons le noyau de $\operatorname{Gal}(k/k_0)\mathbin{\text{---}}G(\mathbb A_{\mathrm{fin}})$, qui contient \unsure{le} noyau $\operatorname{Gal}(k/\bar k_0)$ de $\operatorname{Gal}(k/k_0)\mathbin{\text{---}}\operatorname{Gal}(\unsure{k/\bar k_0|\bar k_0/k_0})$. Ce noyau est envoyé dans $G(\mathbb A_{\unsure{C|c}})$ par $\varphi_{v\infty}$, et l'image est encore dense. \deleted{De façon plus précise,} \inserted{Il me semble que} \deleted{$\varphi_v$ est} \deleted{dominant} $\overline{\varphi_v(\operatorname{Gal}(k/k_0))}$ est un sous (\unsure{grps|gprs}) ouvert de $G(\mathbb A_{\unsure{C|c}})$.
\margin{\unsure{(1)?|( )?}}
\mspage{60}
